% preamble/chapter.title.tex
% 章节扉页设计（蓝色系）：全纸面图片 + 彩条 + 介绍区
% 与常微分方程笔记相同结构，配色改用蓝色系
% =========================================================

\RequirePackage{tikz}
\RequirePackage{graphicx}
\RequirePackage{etoolbox}
\RequirePackage{setspace}

\makeatletter

% -------- 每章图片接口 --------
\newcommand{\chapterimage}{}
\newcommand{\setchapterimage}[1]{\gdef\chapterimage{#1}}

% -------- 章节介绍接口 --------
\newcommand{\chapterintro}{}
\newcommand{\setchapterintro}[1]{\gdef\chapterintro{#1}}

% -------- 版式参数 --------
\newlength{\ChapterTopImageHeight}
\setlength{\ChapterTopImageHeight}{14.2cm}

\newlength{\ChapterRibbonHeight}
\setlength{\ChapterRibbonHeight}{2.35cm}

\newlength{\ChapterIntroTopSep}
\setlength{\ChapterIntroTopSep}{1.15cm}

\newlength{\ChapterIntroWidth}
\setlength{\ChapterIntroWidth}{\dimexpr\textwidth+0.6cm\relax}

\newlength{\ChapterIntroLeft}
\setlength{\ChapterIntroLeft}{\dimexpr\paperwidth/2-\ChapterIntroWidth/2\relax}

\newlength{\ChapterImageMid}
\setlength{\ChapterImageMid}{\dimexpr\ChapterTopImageHeight/2\relax}

\newlength{\ChapterImageShadeTop}
\setlength{\ChapterImageShadeTop}{\dimexpr\ChapterTopImageHeight-1.0cm\relax}

\newlength{\ChapterRibbonMid}
\setlength{\ChapterRibbonMid}{\dimexpr\ChapterTopImageHeight+\ChapterRibbonHeight/2\relax}

\newlength{\ChapterRibbonBottom}
\setlength{\ChapterRibbonBottom}{\dimexpr\ChapterTopImageHeight+\ChapterRibbonHeight\relax}

\newlength{\ChapterIntroTop}
\setlength{\ChapterIntroTop}{\dimexpr\ChapterTopImageHeight+\ChapterRibbonHeight+\ChapterIntroTopSep\relax}

% 章节介绍排版
\newcommand{\ChapterIntroBlock}{%
	\ifdefempty{\chapterintro}{%
		\vspace*{6.4cm}%
	}{%
		{\Large\bfseries\color{BlueMain} 本章导读\par}%
		\vspace{0.55cm}%
		{\setstretch{1.35}\large\itshape\color{BlueDark}\chapterintro\par}%
	}%
}

% 通用章节扉页：#1 = 彩条中的标题文字
\newcommand{\ChapterTitlePage}[1]{%
	\thispagestyle{empty}%
	\vspace*{-\baselineskip}%
	\begin{tikzpicture}[remember picture,overlay]
		% ====== 上部：整页宽图片区 ======
		\ifdefempty{\chapterimage}{%
			\shade[top color=BlueMain!18,bottom color=BlueMain!6]
				(current page.north west) rectangle
				([yshift=-\ChapterTopImageHeight]current page.north east);
			\node[align=center,text=BlueDark,font=\large\itshape]
				at ([yshift=-\ChapterImageMid]current page.north)
				{请在本章前设置章节图片};
		}{%
			\node[anchor=north west,inner sep=0pt,outer sep=0pt]
				at (current page.north west)
				{\includegraphics[
					width=\paperwidth,
					height=\ChapterTopImageHeight
				]{\chapterimage}};
		}%
		% 图片底部渐变
		\shade[top color=black!0,bottom color=black!22]
			([yshift=-\ChapterImageShadeTop]current page.north west) rectangle
			([yshift=-\ChapterTopImageHeight]current page.north east);
		
		% ====== 中部：全宽蓝色彩条 ======
		\fill[BlueMain!94]
			([yshift=-\ChapterTopImageHeight]current page.north west) rectangle
			([yshift=-\ChapterRibbonBottom]current page.north east);
		\fill[GoldMain!75]
			([yshift=-\ChapterTopImageHeight]current page.north west) rectangle
			([xshift=0.18cm,yshift=-\ChapterRibbonBottom]current page.north west);
		\fill[GoldMain!75]
			([xshift=-0.18cm,yshift=-\ChapterTopImageHeight]current page.north east) rectangle
			([yshift=-\ChapterRibbonBottom]current page.north east);
		\node[text=white,font=\bfseries\huge,align=center,text width=\dimexpr\paperwidth-2.2cm\relax]
			at ([yshift=-\ChapterRibbonMid]current page.north)
			{#1};
		\draw[GoldMain!45,line width=0.7pt]
			([yshift=-\ChapterRibbonBottom]current page.north west) --
			([yshift=-\ChapterRibbonBottom]current page.north east);
		
		% ====== 下部：章节介绍区域 ======
		\coordinate (introNW) at
			([xshift=\ChapterIntroLeft,
			  yshift=-\ChapterIntroTop]
			 current page.north west);
		\draw[GoldMain!45,line width=0.8pt] (introNW) -- ++(\ChapterIntroWidth,0);
		\node[anchor=north west,inner sep=0pt,outer sep=0pt,text width=\ChapterIntroWidth,align=left]
			at ([yshift=-0.55cm]introNW)
			{\ChapterIntroBlock};
	\end{tikzpicture}%
	\vspace*{\fill}%
	\clearpage%
}

% numbered chapter 标题
\def\@makechapterhead#1{%
	\if@mainmatter
		\ChapterTitlePage{第\thechapter 章\quad #1}%
	\else
		\vspace*{50pt}%
		{\noindent\huge\bfseries #1\par}%
		\vskip 40pt%
	\fi
}

% unnumbered chapter 标题
\def\@makeschapterhead#1{%
	\if@mainmatter
		\ChapterTitlePage{#1}%
	\else
		\vspace*{50pt}%
		{\noindent\huge\bfseries #1\par}%
		\vskip 40pt%
	\fi
}

\makeatother
